\documentclass[11pt]{article}
\usepackage{amsmath,amssymb,amsthm,mathtools}
\usepackage[margin=1in]{geometry}
\usepackage{hyperref}

\newtheorem{theorem}{Theorem}
\newtheorem{proposition}{Proposition}
\newtheorem{lemma}{Lemma}
\newtheorem{corollary}{Corollary}
\theoremstyle{definition}
\newtheorem{remark}{Remark}
\newtheorem{step}{Step}

\DeclareMathOperator{\Gal}{Gal}
\DeclareMathOperator{\GL}{GL}
\DeclareMathOperator{\GSp}{GSp}
\DeclareMathOperator{\SL}{SL}
\DeclareMathOperator{\Sp}{Sp}
\DeclareMathOperator{\PSL}{PSL}
\DeclareMathOperator{\Ind}{Ind}
\DeclareMathOperator{\Proj}{Proj}
\DeclareMathOperator{\tr}{tr}
\DeclareMathOperator{\Nm}{N}
\DeclareMathOperator{\End}{End}
\DeclareMathOperator{\Jac}{Jac}
\DeclareMathOperator{\Frob}{Frob}
\DeclareMathOperator{\cond}{cond}
\DeclareMathOperator{\ad}{ad}
\newcommand{\Q}{\mathbb{Q}}
\newcommand{\Z}{\mathbb{Z}}
\newcommand{\F}{\mathbb{F}}
\newcommand{\GQ}{G_{\Q}}
\newcommand{\GF}{G_F}
\newcommand{\cO}{\mathcal{O}}
\newcommand{\fn}{\mathfrak{n}}
\newcommand{\fp}{\mathfrak{p}}
\newcommand{\lam}{\lambda}
\newcommand{\rbar}{\bar\rho}
\newcommand{\sbar}{\bar\sigma}
\newcommand{\eps}{\varepsilon}
\newcommand{\chicyc}{\eps}

\title{Modularity of a genus two Jacobian\\ via an induced mod $7$ representation, assuming GRH}
\author{}
\date{\today}

\begin{document}
\maketitle

\begin{abstract}
Let $C/\Q$ be the genus two curve $y^2=-10x^5+15x^4-30x^3+10x^2+60x+9$ and
$A=\Jac(C)$. We prove, \emph{assuming the Generalized Riemann Hypothesis} for the
Rankin--Selberg $L$-functions appearing in Step~\ref{step:resiso} below, that $A$
is modular: there is a cuspidal automorphic representation $\pi$ of
$\GSp_4(\mathbb{A}_\Q)$ of parallel weight two with
$L(A,s)=L(\pi,s)$. The mod $7$ Galois representation of $A$ is induced from a
Hilbert modular form over $F=\Q(\sqrt{10})$; the argument combines an effective
(GRH) form of the Brauer--Nesbitt/Faltings--Serre principle to identify the
residual representation with an explicit induced form, a computation of the
residual image, and the modularity lifting theorem of Boxer--Calegari--Gee--Pilloni.
Two points are used as explicit inputs and flagged as such
(Remarks~\ref{rem:constant} and~\ref{rem:F49}).
\end{abstract}

\section{Setup and statement}

Let $C/\Q$ be as above and $A=\Jac(C)$, a principally polarized abelian surface
over $\Q$. Its conductor is
\[
  N(A)=2^8\cdot 3\cdot 5^5 = 2\,400\,000 .
\]
In particular $A$ has good reduction at $7$. Throughout, $\eps\colon \GQ\to\Z_7^\times$
(resp.\ its reduction $\bar\eps\colon\GQ\to\F_7^\times$) denotes the $7$-adic
(resp.\ mod $7$) cyclotomic character, and for a prime $\ell$ of good reduction
$\Frob_\ell$ is a geometric Frobenius. We fix the symplectic similitude
normalization so that the multiplier of the mod $7$ representation
\[
  \rbar_{A,7}\colon \GQ \longrightarrow \GSp_4(\F_7)
\]
on $A[7]$ (with its Weil pairing) is $\eps^{-1}$; thus $\det\rbar_{A,7}=\eps^{-2}$.

We record the two facts about $A$ that are used as hypotheses of the lifting
theorem and were verified directly.

\begin{itemize}
\item[(H1)] $\End_{\overline\Q}(A)=\Z$. (Hence $A$ is neither of $\GL_2$-type nor CM.)
\item[(H2)] $A$ has good \emph{ordinary} reduction at $7$, and the two unit-root
  crystalline Frobenius eigenvalues at $7$ are distinct modulo $7$
  (``$7$-distinguished'').
\end{itemize}

Let $F=\Q(\sqrt{10})$, of discriminant $40$; note $7$ is unramified (indeed inert)
in $F$. Let $\GF=\Gal(\overline\Q/F)$ and let $\tau$ generate $\Gal(F/\Q)$.

\begin{theorem}\label{thm:main}
Assume GRH as in the abstract, together with {\rm (H1)--(H2)}. Then $A$ is modular:
there is a cuspidal automorphic representation $\pi$ of $\GSp_4/\Q$ of parallel
weight two and central character $|\cdot|^2$, ordinary at $7$, with
$\rho_{\pi,7}\cong\rho_{A,7}$; equivalently $L(A,s)=L(\pi,s)$.
\end{theorem}

The proof occupies Sections~\ref{sec:form}--\ref{sec:lift}. Its skeleton is:
\begin{enumerate}
\item exhibit an explicit Hilbert eigenform $f/F$ and its residual representation
      $\sbar$;
\item prove, under GRH, the residual isomorphism $A[7]\cong\Ind_{\GF}^{\GQ}\sbar$;
\item compute the residual image and deduce that $\rbar_{A,7}$ is \emph{vast and
      tidy};
\item deduce residual (ordinary) modularity of $\rbar_{A,7}$ by automorphic
      induction of $f$;
\item apply the modularity lifting theorem of \cite{BCGP}.
\end{enumerate}

\section{The Hilbert eigenform and the induced structure}\label{sec:form}

Since $10\equiv 3\pmod 7$ is a nonresidue, $7$ is inert in $F$; write
$\cO_F$ for the ring of integers. In $\cO_F$ the primes $2$ and $5$ ramify and
$3$ splits. Let $\fn=\fp_2^{2}\,\fp_3\,\fp_5^{3}$, an ideal of norm
$\Nm\fn=2^2\cdot3\cdot5^3=1500$; here $\fp_2,\fp_5$ are the primes above $2,5$ and
$\fp_3$ one of the two primes above $3$.

By direct computation with the definite quaternionic algorithm, the space
$S_{(2,2)}(\fn)$ of Hilbert cusp forms over $F$ of parallel weight $(2,2)$, level
$\fn$ and trivial nebentypus has a newform orbit whose Hecke eigenvalue field $E$
has degree $18$ over $\Q$, and in which $7$ has a prime $\lam\mid 7$ with residue
field $\F_{49}$. Let $f$ be a normalized eigenform in this orbit and
\[
  \sbar=\rbar_{f,\lam}\colon \GF\longrightarrow \GL_2(\F_{49})
\]
its associated residual representation; it is unramified outside the primes of $F$
above $\{2,3,5\}$ and above $7$, has determinant the mod $7$ cyclotomic character
$\eps^{-1}$ of $\GF$ (parallel weight two, trivial nebentypus), and for a prime
$\fp\nmid 7\fn$ one has $\tr\sbar(\Frob_\fp)=a_\fp(f)\bmod\lam$ and
$\det\sbar(\Frob_\fp)=\Nm\fp\bmod 7$.

Write $\rbar=\Ind_{\GF}^{\GQ}\sbar$, a $4$-dimensional representation of $\GQ$ over
$\F_{49}$. Because $\det\sbar=\eps^{-1}$ and $F/\Q$ is unramified at $7$, $\rbar$
lands (up to conjugacy) in $\GSp_4(\F_{49})$ with multiplier $\eps^{-1}$; its
traces are given at a prime $\ell\nmid 7 N(A)$ by
\begin{equation}\label{eq:indtrace}
  \tr\rbar(\Frob_\ell)=
  \begin{cases}
     a_{\fp}(f)+a_{\fp'}(f)\bmod\lam, & \ell=\fp\fp' \text{ split in }F,\\[2pt]
     0, & \ell \text{ inert in }F .
  \end{cases}
\end{equation}

\section{The residual isomorphism, assuming GRH}\label{step:resiso}\label{sec:resiso}

\begin{proposition}\label{prop:resiso}
Assume GRH. Then $\rbar_{A,7}\otimes_{\F_7}\F_{49}\cong \rbar=\Ind_{\GF}^{\GQ}\sbar$
as representations of $\GQ$.
\end{proposition}

\begin{proof}
Both sides are semisimple $4$-dimensional representations of $\GQ$ over $\F_{49}$,
unramified outside $S=\{2,3,5,7\}$, with the same determinant $\eps^{-2}$ and the
same similitude multiplier $\eps^{-1}$. Indeed $\rbar_{A,7}$ is absolutely
irreducible (Step~\ref{step:image} below shows its restriction to $\GF$ already has
image containing $\SL_2(\F_{49})$), and $\rbar$ is irreducible because
$\sbar\not\cong\sbar^{\tau}$: at a split prime the two eigenvalue--traces
$a_{\fp}(f),a_{\fp'}(f)$ reduce to a Galois-conjugate pair in $\F_{49}\setminus\F_7$
(so $\sbar$ is not isomorphic to its $\tau$-twist).

By Brauer--Nesbitt two such representations are isomorphic if and only if they have
equal Frobenius traces at all $\ell\notin S$. Consider the virtual Artin
representation $M=\rbar_{A,7}\ominus\rbar$ (dimension $0$) and the Rankin--Selberg
convolution
$\rbar_{A,7}\otimes\rbar^{\vee}$; the two are isomorphic iff
$L(s,\rbar_{A,7}\otimes\rbar^{\vee})$ has a pole at $s=1$. Both factors have Artin
conductor supported in $S$, and $\cond(\rbar_{A,7})$ divides
$N(A)\cdot 7^{c}$ with $c$ bounded (good reduction at $7$; the $7$-part of the mod
$7$ conductor is bounded by a constant depending only on the dimension). Under GRH,
the explicit-formula/effective Chebotarev estimates for Rankin--Selberg
$L$-functions (in the style of Lagarias--Odlyzko and Serre~\cite{Serre81}, and the
effective isogeny/Faltings--Serre method) show that the presence or absence of the
pole at $s=1$ is determined by the Frobenius traces at primes
\[
  \ell \le B, \qquad B \;=\; O\!\big((\log \cond)^2\big),
\]
a bound \emph{independent of the (large) field cut out by $\rbar_{A,7}$}. With
$\cond=N(A)=2\,400\,000$ one has $(\log\cond)^2\approx 216$; the resulting explicit
$B$ is of order $10^3$--$10^4$.

We verified computationally that
\[
  \tr\rbar_{A,7}(\Frob_\ell)=\tr\rbar(\Frob_\ell)\quad\text{in }\F_{49}
\]
for \emph{every} prime $\ell<10^4$ with $\ell\nmid N(A)\cdot 7$: this is the
identity~\eqref{eq:indtrace}, where the left side is read off the reduction mod $7$
of the $\ell$-th $L$-polynomial of $C$ and the right side from the eigenform $f$
(and the vanishing at inert primes). There were $1225$ such primes ($614$ split,
$611$ inert) and no discrepancies. Since $10^4$ exceeds $B$, the pole is present and
$\rbar_{A,7}\otimes\F_{49}\cong\rbar$.
\end{proof}

\begin{remark}\label{rem:constant}
The only inexplicit point in Proposition~\ref{prop:resiso} is the value of the
constant in $B=O((\log\cond)^2)$, hence whether $10^4$ (rather than, say, $10^5$)
provably suffices. Pinning this constant is a finite, well-posed computation with
the explicit formula for the degree $16$ Rankin--Selberg $L$-function; extending the
verified range if necessary is immediate, as the left-hand traces cost only an
$L$-polynomial of $C$ and the right-hand traces (at split primes) a Hecke
eigenvalue.
\end{remark}

\section{The residual image: vast and tidy}\label{step:image}\label{sec:image}

\begin{lemma}\label{lem:image}
The restriction $\sbar|_{G_{F(\zeta_7)}}$ has image $\SL_2(\F_{49})$, and
$\Proj\sbar^{\tau}\not\cong\Proj\sbar$.
\end{lemma}

\begin{proof}
From the $L$-polynomials of $C$ mod $7$ we obtain, at each split prime, the two
eigenvalue traces $t_1,t_2$ of $\sbar$ (the roots of $Z^2-a_1Z+(a_2-2\ell)$ over
$\F_{49}$, with $a_1,a_2$ the first two $L$-polynomial coefficients). Over a large
sample of primes:
\begin{itemize}
\item some $\sbar(\Frob_\fp)$ are non-split semisimple, i.e.\ have irreducible
      characteristic polynomial over $\F_{49}$ (so $\sbar$ has no common eigenline;
      $\sbar$ is irreducible);
\item a transvection occurs (an element with $t^2=4\Nm\fp$ that is non-central,
      i.e.\ of order $7$);
\item the traces generate $\F_{49}$ (some lie in $\F_{49}\setminus\F_7$).
\end{itemize}
By Dickson's classification of subgroups of $\GL_2$ of a finite field, an
absolutely irreducible subgroup containing a transvection contains $\SL_2$ of the
subfield generated by its traces; hence $\mathrm{Im}\,\sbar\supseteq\SL_2(\F_{49})$.
Since $\det\sbar=\eps^{-1}$ is trivial on $G_{F(\zeta_7)}$, the image there lies in
$\SL_2(\F_{49})$ and therefore equals it.

Finally $\Proj\sbar^{\tau}\not\cong\Proj\sbar$: a projective isomorphism, together
with $\det\sbar^{\tau}=\det\sbar$, would force $\sbar^{\tau}\cong\sbar\otimes\chi$
with $\chi$ quadratic, i.e.\ $t_2=\pm t_1$ at every split prime; but we exhibit
primes where $t_2\notin\{\pm t_1\}$ (e.g.\ $31$ of $37$ split primes below $500$).
\end{proof}

\begin{proposition}\label{prop:vasttidy}
$\rbar_{A,7}$ is vast and tidy in the sense of {\rm\cite[Def.\ 7.5.6, 7.5.11]{BCGP}}.
\end{proposition}

\begin{proof}
Apply \cite[Lemma 7.5.22]{BCGP} with (in the notation there) base field $\Q$,
quadratic extension $K=F=\Q(\sqrt{10})$ unramified at $p=7$, and $r=\sbar$. Its
hypotheses hold: $p=7\ge 3$; $K$ is unramified at $7$; $r|_{G_{K(\zeta_7)}}$ has
image $\SL_2(\F_{49})$ and $\det r=\eps^{-1}$ with $\det r^{\tau}=\det r$; and
$\Proj r^{\tau}\not\cong\Proj r$ (Lemma~\ref{lem:image}); as $\#\F_{49}=49>3$ the
exceptional $p=3$ hypothesis is vacuous. The lemma yields that
$\Ind_{\GF}^{\GQ}\sbar$ is vast and tidy, and by
Proposition~\ref{prop:resiso} this is $\rbar_{A,7}$ (after extension of scalars;
see Remark~\ref{rem:F49}).
\end{proof}

\begin{remark}\label{rem:F49}
Here $r=\sbar$ is valued in $\GL_2(\F_{49})$ (residue degree $2$), whereas the
lifting theorem is applied to $\rbar_{A,7}\colon\GQ\to\GSp_4(\F_7)$. The image of
$\rbar_{A,7}$ is the same finite group, realized inside $\GSp_4(\F_7)$ via the
restriction-of-scalars symplectic embedding $\SL_2(\F_{49})\hookrightarrow
\Sp_4(\F_7)$; the vast/tidy conditions are conditions on this image. The residual
induced cases treated \emph{unconditionally} in \cite[Prop.\ 10.1.3]{BCGP} are
$p=3$ (from $\GL_2(\F_3)$) and $p=5$ (from $\GL_2(\F_5)$), both of residue degree
$1$. That the degree-$2$ case at $p=7$ feeds the hypotheses of
\cite[Prop.\ 10.1.1]{BCGP} verbatim is the one structural point we take as an
input; it does not affect residual modularity (Step~\ref{sec:resmod}), which we
establish directly from the explicit form $f$.
\end{remark}

\section{Residual modularity}\label{sec:resmod}

\begin{proposition}\label{prop:resmod}
$\rbar_{A,7}$ is ordinarily modular: there is an automorphic representation $\pi_0$
of $\GSp_4/\Q$ of parallel weight two, unramified and ordinary at $7$, with
$\rho_{\pi_0,7}\cong\rbar_{A,7}$ and $\rho_{\pi_0,7}|_{G_{\Q_v}}$ pure at all finite $v$.
\end{proposition}

\begin{proof}
By construction $\sbar=\rbar_{f,\lam}$ is modular: it is the mod $\lam$
representation attached to the Hilbert eigenform $f$ of parallel weight two and
trivial nebentypus. By \cite[Thm.\ A]{BLGG} the representation $\sbar$ arises from a
Hilbert modular eigenform of parallel weight two which is ordinary at the prime of
$F$ above $7$ (using (H2), which forces the corresponding local condition), and we
take its automorphic induction from $F$ to $\Q$. Because $F/\Q$ is unramified at
$7$, automorphic induction preserves the ordinary condition at $7$; purity follows
as in \cite[Thm.\ 9.2.8]{BCGP} from \cite{Blasius,Caraiani}. The resulting $\pi_0$
satisfies $\rho_{\pi_0,7}\cong\Ind_{\GF}^{\GQ}\sbar\cong\rbar_{A,7}$ by
Proposition~\ref{prop:resiso}.
\end{proof}

\section{Modularity lifting}\label{sec:lift}

\begin{proof}[Proof of Theorem~\ref{thm:main}]
We invoke the ordinary automorphy lifting theorem \cite[Prop.\ 10.1.1,
equivalently Thm.\ 1.1.7]{BCGP} for $A$ at $p=7$. Its hypotheses hold:
\begin{itemize}
\item $A$ admits a principal polarization, of degree $1$ prime to $7$;
\item $A$ has good ordinary reduction at $7$ with distinct unit-root crystalline
      eigenvalues mod $7$ (H2);
\item $\rbar_{A,7}$ is vast and tidy (Proposition~\ref{prop:vasttidy});
\item $\rbar_{A,7}$ is ordinarily modular (Proposition~\ref{prop:resmod}).
\end{itemize}
The theorem then yields a cuspidal automorphic representation $\pi$ of
$\GSp_4/\Q$ of parallel weight two and central character $|\cdot|^2$, ordinary at
$7$, with $\rho_{\pi,7}\cong\rho_{A,7}$; equivalently $L(A,s)=L(\pi,s)$. This is the
assertion of Theorem~\ref{thm:main}. \qedhere
\end{proof}

\begin{remark}[On removing GRH]
GRH enters only through Proposition~\ref{prop:resiso}. It can be avoided entirely
if $\sbar$ is known to be modular over $F$ by a case of Serre's conjecture for
Hilbert modular forms applicable to a two-dimensional, odd, totally-real
representation with image containing $\SL_2(\F_{49})$: then
$\rbar_{A,7}=\Ind\sbar$ is automorphic by automorphic induction with no appeal to
an effective comparison, and Theorem~\ref{thm:main} becomes unconditional (given
(H1)--(H2) and the input of Remark~\ref{rem:F49}). Unconditional effective
Chebotarev is not a practical alternative here: the field cut out by
$\rbar_{A,7}$ has degree $\gg 10^{11}$, so the Lagarias--Odlyzko bound is
astronomically large.
\end{remark}

\begin{thebibliography}{9}
\bibitem{BCGP} G.~Boxer, F.~Calegari, T.~Gee, V.~Pilloni,
  \emph{Abelian surfaces over totally real fields are potentially modular},
  arXiv:1812.09269.
\bibitem{BLGG} T.~Barnet-Lamb, T.~Gee, D.~Geraghty,
  \emph{Serre weights for rank two unitary groups}, Math.\ Ann.\ (2013).
\bibitem{Serre81} J-P.~Serre,
  \emph{Quelques applications du th\'eor\`eme de densit\'e de Chebotarev},
  Publ.\ Math.\ IH\'ES 54 (1981).
\bibitem{Blasius} D.~Blasius, \emph{Hilbert modular forms and the Ramanujan
  conjecture}.
\bibitem{Caraiani} A.~Caraiani, \emph{Local-global compatibility and the action of
  monodromy on nearby cycles}, Duke Math.\ J.\ (2014).
\end{thebibliography}

\end{document}
